\documentclass[10pt,a4paper]{article}
\usepackage{amsmath,amssymb,amsthm,graphicx}
\usepackage{amsmath,amssymb,amsthm,hyperref,xfrac}
\providecommand{\skpt}{\smallskip}
\providecommand{\ie}{i.e.\ }
\providecommand{\eg}{e.g.\ }
\providecommand{\resp}{respectively\ }
\providecommand{\Subsection}{\subsection}
\providecommand{\Subsubsection}{\subsubsection}
\usepackage[scr=rsfs,cal=euler]{mathalfa}
\multlinegap0pt
\renewcommand{\theenumi}{\alph{enumi}}
\renewcommand{\theenumii}{\roman{enumii}}
\newenvironment{enumeratei}
{\bgroup\def\theenumi{\roman{enumi}}\def\theenumii{\arabic{enumii}}\begin{enumerate}}
{\end{enumerate}\egroup}
\makeatletter
\let\p@enumii\relax
\makeatother
\newcommand{\ldot}{\mathbin{.}}
\newcommand{\lcdots}{\ldot\cdots\ldot}
\newcommand{\RedefinitSymbole}[1]{%
\expandafter\let\csname old\string#1\endcsname=#1
\let#1=\relax
\newcommand{#1}{\csname old\string#1\endcsname\,}%
}
\RedefinitSymbole{\forall} \RedefinitSymbole{\exists}

\newcommand\mto{\mathchoice{\longmapsto}{\mapsto}{\mapsto}{\mapsto}}
\newcommand{\psfrac}[2]{\sfrac{(#1)}{#2}}
\newcommand{\spfrac}[2]{\sfrac{#1}{(#2)}}

\usepackage[all,cmtip]{xy}\let\labelstyle\textstyle
\usepackage{stmaryrd}
\usepackage{bbm}
\usepackage{mathtools}

\theoremstyle{plain}
\newtheorem{themBJ}{Theorem}
\renewcommand*{\thethemBJ}{\Alph{themBJ}}
\newtheorem{them}{Theorem}[section]
\newtheorem{pro}[them]{Proposition}
\newtheorem{pronotation}[them]{Proposition/Notation}
\newtheorem{prodefinition}[them]{Proposition/Definition}
\newtheorem{lemnotation}[them]{Lemma/Notation}
\newtheorem{cor}[them]{Corollary}
\newtheorem{lem}[them]{Lemma}
\newtheorem*{claim}{Claim}

\theoremstyle{definition}
\newtheorem{vocabulary}[them]{Vocabulary}
\newtheorem{defi}[them]{Definition}
\newtheorem{convention}[them]{Convention}
\newtheorem{notation}[them]{Notation} %?
\newtheorem{remnotation}[them]{Remark/Notation} %?
\newtheorem{definotation}[them]{Definition/Notation}
\newtheorem{remdefi}[them]{Remark/Definition}
\newtheorem{reminder-notation}[them]{Reminder/Notation}
\newtheorem{rem}[them]{Remark}
\newtheorem{ex}[them]{Example}
\newtheorem{reminder}[them]{Reminder}
\newtheorem*{problem*}{Main Problem}
\newtheorem*{note}{Note}
\newtheorem*{fact}{Fact}

\newcommand{\dd}{\mathrm{d}}
\newcommand{\rec}{\mathfrak R}
\newcommand{\id}{\mathrm{Id}}
\newcommand{\join}{\mathrm{Joint}}
\newcommand{\como}{\mathfrak{Q}}
\newcommand{\mq}{\mathfrak{MQ}}
\newcommand{\lev}{\mathrm{Lev}}
\newcommand{\ML}{\mathrm{LK}}
\newcommand{\IT}{\mathrm{IK}}
\newcommand{\disp}{\mathfrak{Disp}}
\newcommand{\argth}{\operatorname{Argth}}
\newcommand{\vol}{\operatorname{Vol}}
\newcommand{\parti}{\operatorname{Part}}
\newcommand{\idk}{\operatorname{id}}

\newcommand{\espace}{E}
\newcommand{\subd}{R}
\newcommand{\Tt}{{\mathcal T}}
\newcommand{\XX}{{\mathcal X}}
\newcommand{\EE}{{\mathcal E}}
\newcommand{\A}{{\mathcal A}}
\newcommand{\Var}{{\mathcal L}}
\newcommand{\Id}{\operatorname{Id}}
\newcommand{\one}{\mathbbm{1}}
\let\mb\relax
\let\mn\relax
\newcommand{\trans}{{}^{\mathrm{t}}\mkern-1mu}
\newcommand{\eqdef}{\coloneqq}
\newcommand{\eqdefup}{\coloneqq}
\newcommand{\espcdot}{\mbox{$\,\cdot\,$}}

\let\eps\varepsilon

\newcommand{\Lg}{\lambda|_{[0,1]}}
\newcommand{\la}{\lambda}
\renewcommand{\P}{\mathbb{P}}
\newcommand{\p}{\mathcal{P}}
\newcommand{\m}{\mathcal{M}}
\newcommand{\ma}{\mathrm{Marg}}
\newcommand{\cC}{\mathcal{C}}\let\C\cC

\newcommand{\R}{\mathbb{R}}
\newcommand{\Q}{\mathbb{Q}}
\newcommand{\N}{\mathbb{N}}
\newcommand{\Z}{\mathbb{Z}}

\newcommand{\AND}{\quad\text{and}\quad}

\newcommand{\E}{\mathbb{E}}
\newcommand{\law}{\operatorname{Law}}
\newcommand{\var}{\operatorname{Var}}
\newcommand{\sto}{\mathrm{sto}}
\newcommand{\stosup}{\operatorname{stosup}}
\newcommand{\stoinf}{\operatorname{stoinf}}
\newcommand{\lcsup}{\operatorname{losup}}
\newcommand{\lcinf}{\operatorname{loinf}}

\newcommand{\leqp}{\preceq_+}
\newcommand{\leqc}{\preceq_{C}}
\newcommand{\leqcs}{\preceq_{C,\sto}}
\newcommand{\leqs}{\preceq_{\text{\rm sto}}}
\newcommand{\geqs}{\succeq_{\text{\rm sto}}}
\newcommand{\e}{\mathrm{e}}

\newcommand{\pr}{\operatorname{proj}}

\newcommand{\pdec}{\mathcal{P}^{\scriptscriptstyle\searrow}}
\newcommand{\mdec}{\mathcal{M}^{\scriptscriptstyle\searrow}}
\newcommand{\leb}{\LL_{[0,1]}}
\newcommand{\op}{\mathopen}
\newcommand{\clo}{\mathclose}
\newcommand{\leqlc}{\preceq_{\rm lo}}
\newcommand{\geqlc}{\succeq_{\rm lo}}
\newcommand{\rsemibracket}{\mathclose{\rfloor\hspace*{-.9ex}\lceil}}
\newcommand{\lsemibracket}{\mathclose{\rfloor\hspace*{-.9ex}\lceil}}
\newcommand{\semi}{\mathclose{\rfloor\hspace*{-.8ex}\lceil}}

\newlength{\arraycolsepsauvegardegenerale}
\setlength{\arraycolsepsauvegardegenerale}{\arraycolsep}
\newlength{\arraycolsepsauvegardetemporaire}

\newlength{\tabcolsepsauvegardegenerale}
\setlength{\tabcolsepsauvegardegenerale}{\tabcolsep}
\newlength{\tabcolsepsauvegardetemporaire}

\newenvironment{nnarray}%
{\setlength{\arraycolsepsauvegardetemporaire}{\arraycolsep}\setlength{\arraycolsep}{0cm}\begin{array}}
{\end{array}\setlength{\arraycolsep}{\arraycolsepsauvegardetemporaire}}

\newenvironment{nntabular}%
{\setlength{\tabcolsepsauvegardetemporaire}{\tabcolsep}\setlength{\tabcolsep}{0cm}\begin{tabular}}
{\end{tabular}\setlength{\tabcolsep}{\tabcolsepsauvegardetemporaire}}
\title{The action of a finite family of increasing kernels preserves the stochastic order and stabilises the decreasing-density class}
\author{The Univalent Foundations Program}
\date{doi:10.5802/jep.177; CC BY-SA 3.0 Unported}
\begin{document}
\maketitle
\setcounter{section}{4}
\setcounter{them}{3}
\noindent\textit{Editorial note.} The counted statement is the author's `lem:mono1` (source lines 1519-1530, Lemma 4.4); the proof environment starts on the next line, so the pairing is by adjacency.\par\medskip
\noindent\textit{Editorial note.} The retained counter values (section 4, theorem offset 3) make the counted lemma print as Lemma 4.4; this article carries no label-number dump, so the number was inferred by counting the theorem-class environments of section 4 and is reported with that caveat.\par\medskip
\noindent\textit{Editorial note.} No retained line refers to a label, citation or figure outside the retained span, so the package declares no external reference; the retained bibliography is bundled because the native preamble of this source is shared with the article's other results.\par\medskip
\noindent\textit{Editorial note.} The retained preamble is the publisher layout adapter followed by the native author macro and theorem block of the pinned source.\par\medskip
\noindent\textit{Editorial scope.} One original result of the pinned publisher article, counted once: monotonicity and stabilisation of the action of a finite kernel family (author label lem:mono1).\par\medskip
\noindent\textit{Allowed standard prerequisites (editorial).}\begin{itemize}
\item Increasing kernel and stochastic order. A kernel is increasing when it preserves the order on measures; join(la; l\_R) is the law obtained by applying the kernel l\_R to la.. It is the first assertion of the statement.
\item Decreasing-density class. mdec(la) is the class of measures with decreasing density with respect to la.. It is the class that the second assertion shows is stabilised.
\item Cumulative distribution functions. The order on measures is tested on cumulative distribution functions, which are affine on the components of A\_r.. It is the argument of the second half of the proof.
\end{itemize}\medskip
Notice that $\ell_{\{r\}}=\ell_r$ and that for any $R$, $\la\ldot\ell_R=\la$. Moreover, $\ell_R$ only depends on $(A_r)_{r\in R}$. The following lemma is particularly simple.

\begin{lem}\label{lem:mono1}
Let $R\subset \R$ be a finite set and $\mu,\,\nu$ be in $\m(\la)$.

\begin{enumerate}
\item
\label{item:mono1a} If $\mu\leqs \nu$ then $\mu \ldot\ell_R\leqs \nu\ldot\ell_R$, \ie $\join(\la; \ell_R)$ has increasing kernel.
\item
\label{item:mono1b} If $\mu$ is in $\mdec(\la)$, then so is $\mu \ldot\ell_R$, and $\mu\leqs \mu\ldot\ell_R$.
\end{enumerate}
\end{lem}

\begin{proof} For $r\in \R$, one easily checks that $\ell_r$ is increasing and stabilizes $\mdec$; \eqref{item:mono1a}~and the first point of \eqref{item:mono1b} follow. If $\mu\in \mdec(\la)$, to see that $\mu\leqs \mu\ldot\ell_R$, look at the cumulative distribution functions. They coincide off the components of $A_r$. Now on each of those, $F_{\mu\ldot\ell_R}$ is affine whereas $F_\mu$ is concave since $\mu$ has decreasing density, so necessarily $F_\mu\geqslant F_{\mu\ldot\ell_R}$.
\end{proof}

\end{document}
